% dhahara.tex
\documentclass[11pt]{book}
\usepackage[utf8]{inputenc}
\usepackage[T1]{fontenc}
\usepackage{geometry}
\geometry{a5paper, margin=1in}
\usepackage{setspace}
\onehalfspacing
\usepackage{microtype}
\usepackage{ebgaramond}
\usepackage{titlesec}
\usepackage{xcolor}
\usepackage{fancyhdr}
\usepackage{array}
\usepackage{tabu}
\usepackage{hyperref}
\hypersetup{colorlinks=true, linkcolor=black, citecolor=black, urlcolor=gray}

\titleformat{\chapter}[hang] 
{\normalfont\LARGE\bfseries}{\thechapter}{1em}{}
\titlespacing*{\chapter}{0pt}{-20pt}{20pt}

\pagestyle{fancy}
\fancyhf{}
\fancyhead[LE,RO]{\thepage}
\fancyhead[RE]{\textit{\leftmark}}
\fancyhead[LO]{\textit{\rightmark}}
\renewcommand{\headrulewidth}{0.4pt}

\newenvironment{lorebox}{\begin{quote}\small\itshape}{\end{quote}}

\begin{document}

\frontmatter
\begin{titlepage}
\centering
\vspace*{2cm}
{\huge\bfseries Dhahara\par}
\vspace{0.5cm}
{\Large Throne of Brass\par}
\vspace{2cm}
{\large A Fate's Edge Expansion\par}
\vspace{1.5cm}
\vfill
\end{titlepage}

\tableofcontents

\mainmatter

\chapter*{Preface: The Two Ledgers of Dhahara}
\addcontentsline{toc}{chapter}{Preface: The Two Ledgers of Dhahara}

Dhahara is not a nation. It is an argument.

In the north, the mountain passes are held by Aeler viceroys who speak a creole of their own tongue, who ride elephants in silver howdahs, who keep fire shrines in the old Dhaharan way. Their ledgers are kept in slate and sealed with keystone wax. They call themselves the **Mountain Satrapies**, and they answer to the High King beneath the peaks – but High Kings are distant, and tax collectors are not.

In the south, the Brass Coast is a patchwork of ports ruled by the **Kahfagian Admiralty**. Here, spice warehouses rise like cathedrals, and the navy patrols the monsoon lanes. The Admiralty styles itself a civilising hand; the Dhaharan call it a velvet glove over an iron hook. The Sidhi – coastal cousins of the Ashaani – are most numerous here, and they remember every slave raid.

Between them lies a no‑man’s land of contested plains, ruined watchtowers, and the great river that some call the boundary and others call a wound. Dhahara has been split for so long that few remember a time when it was whole. The Mountain Viceroy and the Admiral‑Governor meet once a year on a bridge that is neither north nor south, exchange gifts that are also insults, and then return to their ledgers.

Yet Dhahara is far more than the sum of its foreign masters. Beyond the passes and the coast lie independent hill kingdoms, ancient temple‑cities, and riverine principalities where native dynasties still rule. Dhaharan merchants control the inland spice routes; Dhaharan priests maintain the great fire shrines; Dhaharan poets compose epics that the Aeler viceroys themselves commission for their courts. The ledger of Dhahara has two columns, but the ink is mixed with Dhaharan blood, and the interest is paid in Dhaharan grain.

This expansion presents Dhahara as a land of contradictions: monsoon wealth and wartime poverty, elephant cavalry and steam‑driven ropeways, fire‑worshippers and ancestor‑spirit cults. It is also the home of the **Kalrantha** amulet – the green glass shard that a prince stole, that a demon cracked, and that an ancient temple was built to contain. Now the amulet has resurfaced, and the race to return it will tear open every old wound.

When you play in Dhahara, you will bargain with a Viceroy’s chamberlain who weighs your words on golden scales. You will dodge the spies of the Admiralty’s “Trade Protection Bureau”. You will walk the monsoon‑shrouded decks of a Sidhi smuggler’s dhow. And if you are lucky, you will reach the Himadri temple before the amulet’s whispers turn you into the very thing you sought to contain.

---

\chapter{The Land and the Monsoon}

\section{Geography of the Split}

Dhahara is defined by its two halves and the seasonal wind that connects them to the wider world. But those two halves are only the zones of foreign administration. Beyond them lie many other Dhabaras.

\begin{itemize}
    \item \textbf{The Mountain Satrapies (North):} Aeler‑held highlands, terraced for tea and rice. Fortified passes, fire shrines on every peak, and a bureaucracy that measures time in candle‑marks. The Viceroy resides at **Sarvistan**, a city of domed libraries and elephant stables.
    \item \textbf{The Brass Coast (South):} Kahfagian‑administered ports along the Gulf of Oshiira and the eastern ocean. **Port Shed** is the largest, a maze of ropeways, spice auctions, and the Admiralty’s notorious “Maritime Court”. The further east you go, the more the coast blends into Ayokhan trading towns and Sidhi fishing villages.
    \item \textbf{The Scarred Plain:} A buffer of dry grassland, dotted with ruined watchtowers from the last war. Caravans cross it in one day if the weather holds, but bandits and tax inspectors cross it in hours.
    \item \textbf{The Himadri Ridge:} The northernmost spine, where the Aeler passes lead to the high plateau and the secret temple that once held the Kalrantha amulet. The temple is now a ruin, guarded by monks who speak a dialect no one else remembers.
    \item \textbf{The Dhaharan Heartlands:} Beyond the Satrapies and the Admiralty’s reach lie a dozen independent principalities – the hill kingdoms of the eastern ghats, the riverine federation of the middle Ganga, the desert trading cities of the west, and the temple‑states of the southern plateau. Each has its own royal house, its own patron deity, its own courts and coinage. The foreign powers tax what moves between them; they do not rule them.
\end{itemize}

\section{The Monsoon Trade}

The summer monsoon blows from the southwest, bringing ships from Ayokha, Taharka, and the distant eastern isles. The winter monsoon blows from the northeast, carrying Dhaharan silver and Aeler steel back to Kahfagia and Oshiira. This rhythm is the heartbeat of Dhaharan wealth – and the clock that every faction watches.

Use the **monsoon clock [8]** to track the season. At 2, the first ships from Ayokha arrive; markets are flush, prices drop. At 4, the Admiralty raises tariffs. At 6, the northeast winds begin; merchants race to load. At 8, the window closes; any ship still in port is stuck for half a year.

---

\chapter{Peoples of Dhahara}

\section{The Dhaharan Heartlands (Majority Native Cultures)}

\subsection*{Pillars}
Elephant cavalry, temple cities, spice and indigo, ancestor shrines, monsoon festivals.

\subsection*{Voice}
Dhaharan speech is rich in proverbs and layered with honorifics. A farmer greets a merchant with “May the rain find your fields”; a merchant answers a prince with “May your treasury overflow”. Insults are delivered through metaphor, and a well‑turned verse can end a feud.

\subsection*{Regional Variations}
The Dhaharan people are not a single nation but a mosaic of related cultures. The hill kingdoms of the east celebrate the **Feast of the First Rain** with elephant processions and flower‑strewn chariots. The river principalities of the central plains hold **stepwell courts** where disputes are settled at the water’s edge. The desert cities of the west are famed for their **caravanserai poetry**, where traders compete in verse as they haggle over silk and salt. The southern plateau reveres **ancestral spirits** who speak through oracles at mountain shrines.

\subsection*{Etiquette Hooks}
\begin{itemize}
    \item \textbf{Offer a betel leaf} before serious negotiation – refusal is a grave insult.
    \item \textbf{Never step over a threshold with your left foot first} – the left is for the dead.
    \item \textbf{Share the first cup of tea with the eldest present} – even if they are a servant.
\end{itemize}

\subsection*{Strings}
Ancestor tablet (proof of lineage), temple car token (right to a festival seat), weaver’s loom weight (trade credit in the cloth markets).

\section{The Mountain Satrapies (Aeler–Dhaharan)}

\subsection*{Pillars}
Fire shrines, elephant cavalry, tea‑ledgers, mountain hospitality.

\subsection*{Voice}
Satrap officials speak in formal couplets, each line a weight and a counterweight. A joke is a ritual; a threat is a citation of ancient law.

\subsection*{Etiquette Hooks}
\begin{itemize}
    \item \textbf{Present a candle of the right color} – white for petition, red for thanks, black for mourning. (Position +1)
    \item \textbf{Never refuse tea} – the first cup is courtesy, the second is negotiation, the third is dismissal. (Breaking this ends all social rolls for the scene.)
\end{itemize}

\subsection*{Strings}
Fire‑shrine token (right to pray at a public shrine), tea‑ledger page (proof of credit), elephant‑road pass (priority in mountain passes).

\section{The Brass Coast (Kahfagian–Sidhi)}

\subsection*{Pillars}
Spice auctions, ropeways, maritime court, freedman’s pride.

\subsection*{Voice}
Coast merchants speak with Kahfagian cadence but use Dhaharan honorifics. They are direct, suspicious, and quick to invoke the “Maritime Court” as a threat.

\subsection*{Etiquette Hooks}
\begin{itemize}
    \item \textbf{Declare your cargo’s value twice} – once to the clerk, once to the shrine. (The shrine version may be used as a String later.)
    \item \textbf{Never board a ship without a lamp} – even at noon. (Lamp shows you are not a slaver.)
\end{itemize}

\subsection*{Strings}
Spice‑lot stamp (right to trade), ropeway pass (priority loading), freedman’s seal (proof of emancipation from Ashaan).

\section{The Sidhi (Minority from Galanina)}

\subsection*{Pillars}
Sea trade, lighthouse courts, mixed ancestry (Vilikari and Ashaani), fierce independence. The Sidhi are a minority people who trace their origins to the ancient port city of Galanina, far to the west. Generations ago, many fled Ashaani slavery and settled along the Brass Coast. They remain a distinct community – proud, insular, and extraordinarily skilled in navigation and trade.

\subsection*{Voice}
Sidhi speak a creole of Ashaani, Vilikari, and Dhaharan. They are laconic and watchful, with a dark humour about the Admiralty.

\subsection*{Etiquette Hooks}
\begin{itemize}
    \item \textbf{Light a lamp before landing} – Position +1 in Sidhi homes.
    \item \textbf{Never ask about ancestry} – the raids are a fresh wound. (Flips Audience from Warm to Cold.)
\end{itemize}

\subsection*{Strings}
Lighthouse charter, port‑weight (trade advantage), freedman’s seal (honoured in Sidhi courts).

\section{The Imperial Factions (Outsiders with Interests)}

\subsection*{Oshiiran Grain Auditors}
Oshiira sends teams to monitor the grain trade; they are despised but necessary. Their Strings: grain ledger, weigh‑house receipt.

\subsection*{Ashaani “Trade” Envoys}
Veiled Sisters’ agents, officially merchants, actually spying on the Brass Gate. Their Strings: spirit‑lamp, brass seal (counterfeit), a name that is not theirs.

\subsection*{Kahfagian Admiralty Officials}
The colonial administrators, naval officers, and customs clerks who run the Brass Coast. Their Strings: letter of marque, port writ, cipher book.

\subsection*{Ayokhan Merchant Houses}
Wealthy monsoon traders who rent warehouses and hire Sidhi crews. Their Strings: monsoon schedule, temple tithe stamp.

---

\chapter{Patrons of Dhahara}

\section{Regional Faces of Known Patrons}

\subsection{Mykkiel – The Judge of the Peak}
In the Mountain Satrapies, Mykkiel is not a distant arbiter but a local magistrate. His temples are audience halls, his rites are public hearings. He demands that every oath be witnessed by two elders. The Veiled Sisters fear him because he has already tried them in absentia.

\textbf{Regional Rite – The Unsealed Testimony:} Once per session, compel a spirit to speak truth for one question; the spirit’s former master learns your face. (Mark +2 Obligation)

\subsection{Aveh – The Dust‑Rider of the Scarred Plain}
Here, Aveh is a djinn – a spirit of the sandstorm, the liberator of slaves, the patron of those who flee without paying. She appears as a rider without a face, offering an exit from any prison or contract. Her price is always the same: you must never close a door behind you.

\textbf{Etiquette:} When you flee, leave a coin at the threshold. (The GM may convert the first SB of pursuit into a “missed footing” complication instead of harm.)

\textbf{Rite – The Unlatched Chain:} Break one physical bond (shackle, rope, locked door) as a free action. (Mark +1 Obligation; the broken object may later reform.)

\subsection{Ikasha – The Lotus in the Well (Sidhi variant)}
Among the Sidhi, Ikasha is the dream that rises from dark water. She offers refuge to escaped slaves and spirits. Her shrines are wells with a single lamp floating on the surface.

\textbf{Rite – The Black Water:} Once per journey, you may hide one object or person in a Sidhi well for a season; no one will look there. (Mark +1 Corruption; the well may ask for a story in return.)

\section{New Patrons of Dhahara}
\emph{The following patrons are worshipped across the Dhaharan heartlands, from the hill kingdoms to the desert cities. They have no single temple; their shrines are everywhere – at fords, in marketplaces, under ancient banyan trees.}

\subsection{Surya – The Unblinking Ledger}
\paragraph{Domain:} The sun, audits, the moment of truth.
\paragraph{Lore:} Surya is a spirit of the high passes, where the sun bakes the stone and shadows are short. He does not speak; he witnesses. His followers are tax collectors, truth‑seekers, and those who fear the dark.
\paragraph{Favored Deeds:} Exposing a liar, paying a debt at noon, climbing to a shrine.
\paragraph{Taboos:} Lying in direct sunlight, hiding a tax receipt, eating before noon.
\paragraph{Low Rite – The Shadowless Hour (4 XP):} For one scene, you cannot be deceived by visual illusions. (Mark +1 Fatigue)
\paragraph{Standard Rite – Surya’s Gaze (8 XP):} Target must tell one truth about themselves or suffer –1 die to all rolls until dawn. (Mark +2 Obligation)
\paragraph{High Rite – The Unforgiving Ledger (14 XP):} Name a debt; the debtor cannot sleep until they pay. (Mark +3 Corruption; the debt may be spiritual, not monetary.)

\subsection{Chandra – The Well That Reflects}
\paragraph{Domain:} Moon, hidden water, second chances.
\paragraph{Lore:} Chandra is the spirit of the well where a Sidhi child hid from slavers. She is kind, cautious, and forgetful – but she remembers cruelty. Her followers are smugglers, midwives, and the keepers of way‑stations.
\paragraph{Favored Deeds:} Hiding a fugitive, drawing water at night, whispering a secret to the moon.
\paragraph{Taboos:} Using a well to drown evidence, silencing a child’s cry, reflecting a lie in water.
\paragraph{Low Rite – The Moon Pool (4 XP):} Look into any reflective surface; see the nearest safe hiding place. (Mark +1 Fatigue)
\paragraph{Standard Rite – Chandra’s Cup (8 XP):} A waterskin becomes “memory” – anyone who drinks from it recalls a forgotten fact. (Mark +2 Obligation)
\paragraph{High Rite – The Second Well (14 XP):} Create a temporary well that lasts one night. The water heals 1 Harm but imposes a geas. (Mark +3 Obligation)

\subsection{Agni – The Three Fires}
\paragraph{Domain:} Hearth, forge, funeral pyre.
\paragraph{Lore:} Agni is three spirits in one – the cook fire, the smith’s flame, and the cremation torch. He is the patron of cooks, blacksmiths, and undertakers. He has no opinion on war; he cares only about heat.
\paragraph{Favored Deeds:} Keeping a kitchen fed, forging a tool for a worthy cause, performing a proper funeral.
\paragraph{Taboos:} Letting a hearth go cold in winter, using a forge for weapons of assassination, burning a body without witness.
\paragraph{Low Rite – The Smell of Smoke (4 XP):} For one scene, you know the location of the nearest fire. (Mark +1 Fatigue)
\paragraph{Standard Rite – Agni’s Seal (8 XP):} Touch a metal object; it becomes fire‑resistant for one day. (Mark +2 Obligation)
\paragraph{High Rite – The Three Fires (14 XP):} Once per campaign, when you die, you may be cremated immediately. Your ashes grant a boon to your heirs (GM determines). (Mark +3 Corruption; you cannot be resurrected.)

\subsection{Vayu – The Wind That Writes}
\paragraph{Domain:} Monsoon, messages, the speed of news.
\paragraph{Lore:} Vayu is a spirit of the monsoon wind, who carries rumours faster than horses. He is fickle, generous, and easily bored. His followers are messengers, spies, and sailors.
\paragraph{Favored Deeds:} Delivering a message through a storm, reading a wind‑chart correctly, warning a village of a flood.
\paragraph{Taboos:} Burning a letter without reading it, silencing a messenger, hoarding wind‑charms.
\paragraph{Low Rite – The Written Breath (4 XP):} Whisper a message into a leaf; the wind will carry it to the nearest temple of Vayu. (Mark +1 Fatigue)
\paragraph{Standard Rite – Vayu’s Chart (8 XP):} For one leg, your ship or caravan gains +1 Speed. (Mark +2 Obligation)
\paragraph{High Rite – The Monsoon’s Secret (14 XP):} Once per season, you may predict a storm’s path with perfect accuracy. (Mark +3 Corruption; the prediction must be shared freely.)

\subsection{Prithvi – The Elephant That Remembers}
\paragraph{Domain:} Ancestors, memory, the weight of tradition.
\paragraph{Lore:} Prithvi is a spirit of the elephant herds that once roamed the Dhaharan plains. He is slow, patient, and never forgets a face. His followers are cavalrymen, archivists, and those who keep family ledgers.
\paragraph{Favored Deeds:} Recording a family history, protecting a herd, eating an elephant’s favourite fruit.
\paragraph{Taboos:} Selling elephant ivory, forgetting a debt to a relative, riding a pregnant elephant.
\paragraph{Low Rite – The Stomp of Memory (4 XP):} For one scene, you gain +1 die to recall any fact about Dhaharan history. (Mark +1 Fatigue)
\paragraph{Standard Rite – Prithvi’s Seal (8 XP):} Touch an object; you know the name of its maker. (Mark +2 Obligation)
\paragraph{High Rite – The Herd’s Witness (14 XP):} Summon an ancestor spirit for one question; the spirit may lie if it is protecting a secret. (Mark +3 Corruption)

\subsection{Ganga – The River That Divides}
\paragraph{Domain:} Boundaries, fords, the space between north and south.
\paragraph{Lore:} Ganga is the spirit of the great river that splits Dhahara. She is neutral, patient, and merciless to those who try to cross her without a proper ferry. Her followers are ferrymen, bridge‑wardens, and diplomats.
\paragraph{Favored Deeds:} Building a safe crossing, rescuing a drowning person, cleaning a riverbank.
\paragraph{Taboos:} Polluting the river, charging a corpse for passage, cutting a bridge without warning.
\paragraph{Low Rite – The Ford’s Price (4 XP):} Know the safest crossing point on any river. (Mark +1 Fatigue)
\paragraph{Standard Rite – Ganga’s Witness (8 XP):} Any oath sworn on the river is witnessed by Ganga; breaking it imposes –2 dice to all social rolls until atonement. (Mark +2 Obligation)
\paragraph{High Rite – The Divided Waters (14 XP):} Once per season, you may part a narrow river for one minute. (Mark +3 Corruption; the river may ask for a sacrifice in return.)

\subsection{Durga – The Silence of the Blade}
\paragraph{Domain:} Justice, execution, the end of argument.
\paragraph{Lore:} Durga is a spirit of the old Dhaharan courts – the executioner’s sword that fell on oath‑breakers. She is cold, final, and without mercy. Her followers are judges, assassins, and those who have been wronged beyond forgiveness.
\paragraph{Favored Deeds:} Carrying out a lawful execution, protecting a witness, destroying a false idol.
\paragraph{Taboos:} Drawing the blade without a witness, killing a prisoner who has surrendered, forging a death warrant.
\paragraph{Low Rite – The Unsheathed (4 XP):} For one scene, your melee attacks gain +1 die against oath‑breakers. (Mark +1 Fatigue)
\paragraph{Standard Rite – Durga’s Judgement (8 XP):} Name a target; for one day, any lie they tell will be exposed by a chill wind. (Mark +2 Corruption)
\paragraph{High Rite – The Final Argument (14 XP):} Once per campaign, you may kill a mortal oath‑breaker with a single blow – but you must also destroy their name from all records. (Mark +3 Corruption; you become a fugitive from the victim’s family.)

\subsection{Lakshmi – The Random Coin}
\paragraph{Domain:} Luck, windfall, the unexpected profit.
\paragraph{Lore:} Lakshmi is a spirit of the marketplace – the coin that is found, the trade that is offered by a stranger, the spice shipment that arrives early. She is capricious, kind to beggars, and cruel to misers.
\paragraph{Favored Deeds:} Giving alms, playing dice for a friend’s debt, sharing a surprise profit.
\paragraph{Taboos:} Hoarding a found coin, rigging a game of chance, refusing to gamble when the stakes are low.
\paragraph{Low Rite – The Found Penny (4 XP):} You find a single coin on the ground. It is always exactly the amount you need for a small purchase. (Mark +1 Fatigue; the coin vanishes after use.)
\paragraph{Standard Rite – Lakshmi’s Smile (8 XP):} Once per session, you may reroll a failed trade or negotiation roll. (Mark +2 Obligation)
\paragraph{High Rite – The Floating Market (14 XP):} For one hour, you may buy or sell any mundane item in any settlement, regardless of inventory. (Mark +3 Corruption; prices are set by the GM, not market logic.)

---

\chapter{The Kalrantha Amulet}

\section{What It Is}
A small, tear‑shaped amulet of green glass, set in a brass frame that never tarnishes. Warm to the touch when someone nearby is about to decide.

\section{Its True History (GM Only)}
The amulet was created by the priests of the Himadri temple as a trap for a demon of indecision. They bound the demon into the glass, and the temple’s wards were designed to keep it dormant. A Dhaharan prince, in his arrogance, stole it centuries ago. The demon whispered temptations; the prince conquered a kingdom, then lost it. The amulet was lost in the desert. The demon is not dead – only sleeping. The temple is now a ruin, its wards cracked. The amulet wants to be used again. It wants to *warp fate* for its wielder.

\section{Game Mechanics (As Known to Players)}
The amulet appears to do very little. Once per session, the wearer may reroll a single die. That is all. Scholars debate its true power. Some say it cursed the prince; others say it was cursed *by* him. The truth is that the amulet is a liability. Every time it is used, the demon stirs. The GM secretly tracks a **Fate Corruption** clock [6]. Each use ticks it. When the clock fills, the amulet manifests a major twist – a betrayal, a sudden storm, a forgotten debt that must be paid in blood. The players will not know the clock exists, but they will feel the pattern.

\section{The Showpiece Adventure: Return of the Kalrantha}

\subsection{The Setup}
An archaeologist – a Dhaharan scholar named **Vikram of the Broken Wheel** – has found the amulet in the Ashaani desert. He knows its history. He wants to return it to the Himadri temple before the demon wakes. He is dead within the first scene, poisoned by a Veiled Sister’s agent. The party is hired by his Tulkani assistant to finish the job.

\subsection{The Race}
Four factions want the amulet:
\begin{itemize}
    \item \textbf{Ashaani (Veiled Sisters):} To use it as a weapon against their enemies. Their agent is a charming merchant named **Nadia** who offers a fortune.
    \item \textbf{Oshiiran Grain Auditors:} They fear the amulet will destabilise the grain trade. Their agent is a cold woman named **Mei‑Lin** who offers safe passage.
    \item \textbf{Kahfagian Admiralty:} They want to study it, possibly to use against Oshiira. Their agent is a cynical privateer named **Captain Red‑Hook** who offers a ship.
    \item \textbf{Dhaharan Monk (the True Believer):} A silent figure in burnt orange robes, **Brother Ananda**, who wants only to seal the amulet. He has no money, but he knows the temple’s wards.
\end{itemize}
The party must travel from Ashaan → a Sidhi port (Port Shed) → north through Dhahara (the Scarred Plain) → the Himadri temple. Each leg has a betrayal, a deal, and a choice of which faction to trust.

\subsection{Key Scenes}
\begin{enumerate}
    \item \textbf{The Poisoned Scholar:} Vikram dies in a crowded market. His assistant gives the party the amulet and a map. The murderer is watching.
    \item \textbf{The Three Offers:} In Port Shed, each faction makes a pitch. Brother Ananda offers only the truth: the amulet will kill its user. The others offer coin, safe passage, or a ship.
    \item \textbf{The Scarred Plain:} The party is betrayed by the faction they trusted most (or by a double agent). A chase across the plain, with a sandstorm (courtesy of Aveh) as cover.
    \item \textbf{The Mountain Pass:} The Aeler viceroy demands a toll. The party must pay in story or silver. The amulet trembles.
    \item \textbf{The Himadri Temple:} The wards are failing. Brother Ananda performs a ritual that requires the amulet to be placed on an altar. The other factions attack. The demon wakes. The party must hold the line for five rounds while the ritual completes.
    \item \textbf{The Choice:} In the final moment, the amulet offers the wielder a wish – anything they desire. Accepting it will release the demon. Refusing it will seal it again, but the demon will whisper every night for the rest of the wielder’s life. There is no third option.
\end{enumerate}

\subsection{Twists}
\begin{itemize}
    \item Vikram’s assistant is actually a Black Hand agent (Ashaan) in disguise; she has been guiding the party to remove the other factions.
    \item The Oshiiran agent is not an auditor but a Narrow Captain who wants to start a war.
    \item The Kahfagian captain is a Sidhi who lost his family to the Admiralty; he plans to sell the amulet to Ayokha.
    \item Brother Ananda is not a monk; he is the last priest of the original temple, and he intends to destroy the amulet – not just seal it – which would kill him.
\end{itemize}

\subsection{Resolution}
If the party seals the amulet, they gain a permanent boon: **The Demon’s Debt** – once per session, they may reroll a failed die, but the GM gains 1 SB. (The demon still stirs.) If they destroy it, Brother Ananda dies, and the party gains a curse: **The Unsealed Vow** – the demon will hunt them in their dreams (roll a Dread Clock each session). If they use the wish, the demon is freed; the campaign shifts to a new arc: hunting the demon.

---

\chapter{Adventure Framework}

\section{Adventure Seeds (Ten Drop‑Ins)}
\begin{enumerate}
    \item \textbf{The Elephant’s Ledger:} A Mountain Satrapy official is selling stolen elephant‑road passes. The party must audit his accounts without revealing their true purpose.
    \item \textbf{The Smell of Monsoon:} A Sidhi crew is accused of smuggling spice; the real cargo is escaped slaves. The party must get them to a safe house.
    \item \textbf{The Forged Seal of the Admiralty:} A forger is creating false letters of marque. The Admiralty hires the party to find him; the forger is a former slave.
    \item \textbf{The Fire That Wouldn’t Die:} A fire shrine in the mountains has been burning for a year without fuel. The priest is drawing on a deep vein of magic – and something is answering.
    \item \textbf{The River Court:} Ganga’s spirit is offended; the river has risen and flooded a village. The party must arbitrate between the north and south banks.
    \item \textbf{The Random Coin:} Lakshmi appears as a beggar and offers the party a single coin. If they spend it wisely, they gain a fortune; if they hoard it, they lose everything.
    \item \textbf{The Silence of the Blade:} A judge is executing the wrong man. The party must prove the real killer’s guilt – or perform the execution themselves.
    \item \textbf{The Three Fires:} A smith has forged a sword that burns with Agni’s flame. He wants to sell it to the Viceroy; the Admiralty wants to steal it; the party is caught in the middle.
    \item \textbf{The Wind That Writes:} A monsoon message is intercepted. It contains a treaty that would end the cold war. The party must deliver it before the signatories are killed.
    \item \textbf{The Well That Reflects:} A Sidhi girl is hiding in a well; the Black Hand is searching for her. The party must get her to the coast safely.
\end{enumerate}

\section{GM Toolkit}

\subsection{Weather and Hazards (d12)}
\begin{enumerate}
    \item Monsoon squall – all sails Desperate for a leg.
    \item Dust storm – Vision –1 Position.
    \item Landslide – mountain pass closed; detour adds 2 segments.
    \item Elephant herd crossing – delay 1 day; harm if startled.
    \item River flood – ferry destroyed; lose 1 Supply.
    \item Earthquake – temple ruins shift; new paths open.
    \item Smoke from a forest fire – test Endurance or mark Fatigue.
    \item Perfect calm – no wind; ships drift.
    \item Meteor shower – omens; the GM gains 1 SB.
    \item Ashfall – from distant volcano; all actions Desperate.
    \item Leeches – in the Scarred Plain’s wet season; lose 1 Supply to treat.
    \item The Ninth Hour – sunset twice as long; the GM may ask a player to describe an omen.
\end{enumerate}

\subsection{Factions of Dhahara}
\begin{tabu} to \linewidth {|X|c|c|c|X|}
\hline
\textbf{Faction} & \textbf{Influence} & \textbf{Stability} & \textbf{Crisis} & \textbf{Goal} \\
\hline
Mountain Satrapy & 6 & 5 & 4 & Keep the passes closed to the Admiralty \\
Brass Coast Admiralty & 6 & 4 & 6 & Monopolise spice trade \\
Sidhi Port Guilds & 4 & 3 & 7 & Free the coast from Ashaan \\
Ashaani Agents & 4 & 3 & 5 & Recover Kalrantha \\
Oshiiran Auditors & 5 & 5 & 3 & Stabilise grain prices \\
Ayokhan Merchant Houses & 5 & 5 & 4 & Profit from both sides \\
\hline
\end{tabu}

\chapter*{Colophon}
\addcontentsline{toc}{chapter}{Colophon}

\vspace{2\baselineskip}
\begin{lorebox}
Dhahara is a ledger with two columns. The Mountain writes in slate; the Coast writes in ink. Neither column ever totals. The amulet Kalrantha is the missing line – the error that neither side can account for. It was stolen from a temple that should not have been built, by a prince who should not have been born, and it has been running through the balance sheets of history ever since.

When you return it, you will not balance the ledger. You will only stop it from bleeding.

— \textit{Dusana of the Raven Road, at Port Shed, during a monsoon that refused to end. The ink ran in rivulets.}
\end{lorebox}

\end{document}